\documentclass[10pt, a4paper, twocolumn]{article}

% --- 基本設定 ---
\usepackage{geometry}
\geometry{a4paper, top=1.8cm, bottom=0.5cm, left=0.5cm, right=0.5cm}

% --- 中文支援 ---
\usepackage{xeCJK}
\setCJKmainfont{Noto Serif CJK TC} % 可替換為其他字型，如 Microsoft JhengHei 或 PingFang TC

% --- 數學與符號 ---
\usepackage{amsmath}

% 支援刪除線
\usepackage[normalem]{ulem}

% 段落縮排
\usepackage{indentfirst}
\setlength{\parindent}{2em} 
% 行距調整
\usepackage{setspace}
\usepackage{enumitem}

% --- 預設內文樣式 ---
\setstretch{1.3}
\setlength{\parskip}{0.6em}

% --- 定義還原樣式（單行距 + 無段距） ---
\newcommand{\normaltextspacing}{%
	\setstretch{1.0}%
	\setlength{\parskip}{0pt}%
}
\newcommand{\resumetextspacing}{%
	\setstretch{1.3}%
	\setlength{\parskip}{0.6em}%
}

% --- 進入這些區塊時暫時關閉文字樣式 ---
\BeforeBeginEnvironment{itemize}{\normaltextspacing}
\AfterEndEnvironment{itemize}{\resumetextspacing}

\BeforeBeginEnvironment{enumerate}{\normaltextspacing}
\AfterEndEnvironment{enumerate}{\resumetextspacing}

\BeforeBeginEnvironment{minted}{\normaltextspacing}
\AfterEndEnvironment{minted}{\resumetextspacing}

\BeforeBeginEnvironment{minted}{\vspace{-0.5em}}
\AfterEndEnvironment{minted}{\vspace{-0.8em}}

% tight itemize
\SetEnumitemKey{tight_item}{
    nosep,
    leftmargin=*,
    after={\vspace{-0.8em}}
}

% 表格行距
\renewcommand{\arraystretch}{1.5}

% 自訂 quote
\usepackage[most]{tcolorbox}
\usepackage{xcolor}

\renewenvironment{quote}
{\begin{tcolorbox}[
		colback=white,
		colframe=gray,
		left=4mm,
		right=0mm,
		top=2mm,
		bottom=2mm,
		boxrule=-1pt,
		enhanced,
		borderline west={1.5pt}{0pt}{gray},
		sharp corners,
		before skip=8pt, after skip=8pt,
		boxsep=2pt,
		breakable, 
		fontupper=\color{gray}
		]}
	{\end{tcolorbox}}

\ifdefined\isdraft
    \usepackage[verbatim]{minted}  % 關掉 highlighting 快速編譯
\else
    % \usepackage[verbatim]{minted}
    \usepackage{minted}
\fi

\setminted{
	fontsize=\small,
	frame=lines,   % 加上外框
	breaklines     % 自動換行
}

% --- 圖片與超連結 ---
\usepackage{graphicx}
\usepackage[hidelinks]{hyperref}
\usepackage{float} % 用於 [H] 強制圖片位置

% --- 標題與作者 ---
\title{TLE 的毒瘤高科技}
\author{暴力又被TLE}
\date{\today}

% --- 頁首頁尾設定 ---
\usepackage{fancyhdr}
\pagestyle{fancy}
\fancyhf{} % 清除所有預設的頁首頁尾

% 使用 \makeatletter 以便存取 \@title 和 \@author 內部變數
\makeatletter
\lhead{\@title}
\rhead{\thepage}
\makeatother

\setlength{\headsep}{12pt} % 減少頁首與內文的間隔
\renewcommand{\headrulewidth}{0.4pt} % 頁首底下的分隔線粗細
\renewcommand{\footrulewidth}{0pt} % 移除頁尾分隔線 (若有)

% --- 縮小 section ---
\usepackage{titlesec}
% 格式: \titleformat{命令}[型態]{格式}{標號}{標號與文字間距}{標題內容前綴}
\titleformat{\section}
  {\normalfont\large\bfseries}   % 樣式 (這裡把 section 改小成 large)
  {\thesection}                  % 顯示章節號
  {1em}                          % 章節號和文字之間的距離
  {}                             % 標題內容的前綴 (這裡留空)
\titleformat{\subsection}
  {\normalfont\normalsize\bfseries}
  {\thesection}
  {1em}
  {}

% 格式: \titlespacing{命令}{左邊縮排}{與前段距離}{與後段距離}
\titlespacing{\section}
  {0pt}      % 左縮排
  {3pt plus 2pt minus 2pt}    % 前段距離 (預設大約 12pt~20pt)
  {3pt plus 2pt minus 2pt}    % 後段距離 (預設大約 6pt~12pt)
\titlespacing{\subsection}
  {0pt}
  {6pt plus 2pt minus 2pt}
  {3pt plus 2pt minus 2pt}

\usepackage[titles]{tocloft}
\setlength{\cftbeforesecskip}{1.2em}

\usepackage{multicol}

\renewcommand{\contentsname}{目錄}

\begin{document}

\begin{multicols}{2}
  \begin{spacing}{0.2}
    \tableofcontents
  \end{spacing}
\end{multicols}

\newpage
	
\section{C++ Template}

\begin{minted}{cpp}
#include <bits/stdc++.h>
#define IO cin.tie(0)->sync_with_stdio(0);
#define R(x) istringstream stm(x); cin.rdbuf(stm.rdbuf());
#define M(x, v) memset(x, v, sizeof(x))
#define REP(i, s, t) for (int i = (s); i < (t); ++i)
#define in >> 
#define ot << 
#define se ot ' '
#define sep se ot 
#define nl ot '\n'
#define all(x) x.begin(), x.end()
using namespace std;
template <typename T>
using vec = vector<T>;
using ll = long long;
void solve();
string T = R"()";

int main() {
  IO;
  int t = 1;

  R(T);
  cin in t;

  while (t--) solve();
  return 0;
}

const int N = 1e5, mod = 1e9 + 7;
const ll inf = 0x3f3f3f3f'3f3f3f3f;

void solve() {
  
}
\end{minted}

\subsection*{Compile}

\begin{minted}{text}
g++ -std=c++20 -Wall -pipe -g3 -O0 test.cpp -o test
g++ -std=c++20 -Wall -pipe -O2 test.cpp -o test
... && (echo [done] && test.exe & cd .) || echo [fail]
...; if ($?) { echo "[done]"; ./test.exe; } else { echo "[fail]"; }
if ...; then echo "[done]"; ./test; else echo "[fail]"; fi
\end{minted}

\subsection*{Pragma}

\begin{minted}{cpp}
#pragma GCC optimize("O3,unroll-loops,no-stack-protector")
#pragma GCC target("avx2,bmi,bmi2,lzcnt,popcnt")
\end{minted}

\subsection*{fread}

\begin{minted}{cpp}
#define gc() (p1==p2&&(p2=(p1=buf)+fread(buf,1,1<<25,stdin),p1==p2) ? EOF : *p1++)
char buf[1<<25],*p1,*p2;
int gi(){char c=gc();int x=c!='-'?c^'0':0,flag=c=='-'; for(;'0'<=(c=gc())&&c<='9';x+=c^'0')x*=10; return flag?-x:x;}
unsigned gu(){char c;unsigned x=0; for(;'0'<=(c=gc())&&c<='9';x+=c^'0')x*=10; return x;}
\end{minted}

\section{PyPy recursion}

\begin{minted}{python}
from sys import setrecursionlimit
setrecursionlimit(10**7)
import pypyjit
pypyjit.set_param('max_unroll_recursion=-1')
\end{minted}

\section{對拍}

\begin{minted}{python}
import subprocess
from itertools import count
from difflib import ndiff

def compile(name):
  subprocess.run(["g++", "-O2", f"{name}.cpp", "-o", f"{name}.exe"])
  print(f"{name} compiled")

compile("test")

gen = ["python", "gen.py"]
ac = ["./test.exe"]
bf = ["python", "bf.py"]

for it in count():
  tc = subprocess.check_output(gen, text=True)
  x = [line.rstrip() for line in subprocess.check_output(ac, input=tc, text=True).splitlines()]
  y = [line.rstrip() for line in subprocess.check_output(bf, input=tc, text=True).splitlines()]
  if x != y:
    print("-" * 30)
    print(tc)
    print("-" * 30)
    print("\n".join(ndiff(x, y)))
    break
  if it % 100 == 0: print(f"pass {it = }")
\end{minted}

\section{BIT}

\begin{minted}{cpp}
struct BIT {
  int n;  // n != 0
  vec<ll> bit;
  BIT(int nn) : n(nn), bit(n + 1) {}
  void add(int i, ll v) {  // l[i] += v
    for (++i; i <= n; i += i & -i) bit[i] += v;
  }
  ll query(int s, int t) {  // sum(l[s:t])
    ll res = 0;
    for (; t > s; t &= t - 1) res += bit[t];
    for (; s > t; s &= s - 1) res -= bit[s];
    return res;
  }
  int kth(ll k) {  // sum(l[:i]) <= k < sum(l[:i+1])
    int i = 0;  // 0 <= i <= n
    for (int b = 1 << __lg(n); b; b >>= 1) {
      if (i + b <= n and bit[i + b] <= k) {
        k -= bit[i += b];
      }
    }
    return i;
  }
};
\end{minted}

\section{PBDS}

\begin{itemize}[tight_item]
    \item 若要 \texttt{multiset} 則用 \texttt{pair\{value, idx++\}}
\end{itemize}

\begin{minted}{cpp}
#include <ext/pb_ds/assoc_container.hpp>
#include <ext/pb_ds/tree_policy.hpp>
using namespace __gnu_pbds;
tree<pii, null_type, less<pii>, rb_tree_tag, tree_order_statistics_node_update> u;
u.order_of_key(x)  // count < x
u.find_by_order(k)  // kth 0-based
\end{minted}

\section{稀疏表}

\begin{minted}{cpp}
template <typename T, auto f>
struct ST {
  static constexpr int M = __lg(N) + 1;
  T st[M][N];

  void build(T *l, int n) {
    const int lim = n + 1;
    int bl = __lg(n) + 1;

    REP(i, 0, n) st[0][i] = l[i];
    int b = 1;
    REP(r, 1, bl) {
      REP(i, 0, n - b) {
        st[r][i] = f(st[r - 1][i], st[r - 1][i + b]);
      }
      b <<= 1;
    }
  }
  T query(int s, int t) {
    int r = __lg(t - s);
    return f(st[r][s], st[r][t - (1 << r)]);
  }
};
// need C++ 20
ST<int, [](auto a, auto b) {return max(a, b);}> st;
\end{minted}

\section{FHQ-Treap}

\begin{itemize}[tight_item]
    \item TLE/MLE：檢查 \texttt{node.siz}、\texttt{siz(o)}、\texttt{rng} 是否 \texttt{long long}
    \item WA：\texttt{update} 前需要 \texttt{clone}
    \item WA：如果 \texttt{merge} 前沒有 \texttt{split} 則需要 \texttt{clone}
    \item WA：所有地方都加上 \texttt{push}
\end{itemize}

\begin{minted}{cpp}
mt19937 rng(114514);
// mt19937_64 rng(1145141919810);
struct node {
  int val, siz;
  node *le, *ri;
  node(int v): val(v), siz(1), le(0), ri(0) {}
  void push() {}
  void pull() {
    siz = 1;
    if (le) siz += le->siz;
    if (ri) siz += ri->siz;
  }
  node *clone() {
    // clone before update!!!
    return this;
    // return new node(*this);  // persistent
  }
};
auto siz(node *x) {
  return x ? x->siz : 0;
}
node *merge(node *a, node *b) {
  if (not a or not b) return a ? a : b;
  // if no split before merge, need clone
  if (rng() % (siz(a) + siz(b)) < siz(a)) {
    a->push(), a->ri = merge(a->ri, b), a->pull();
    return a;
  } else {
    b->push(), b->le = merge(a, b->le), b->pull();
    return b;
  }
}
void split(node *o, node *&a, node *&b, int val) {
  if (not o) {
    a = b = 0; return;
  }
  o = o->clone();
  o->push();
  if (o->val < val) {
    a = o, split(o->ri, a->ri, b, val);
  } else {
    b = o, split(o->le, a, b->le, val);
  }
  o->pull();
}
void split2(node *o, node *&a, node *&b, int k) {
  if (not o) {
    a = b = 0; return;
  }
  o = o->clone();
  o->push();
  if (siz(o->le) < k) {
    a = o, split2(o->ri, a->ri, b, k - siz(o->le) - 1);
  } else {
    b = o, split2(o->le, a, b->le, k);
  }
  o->pull();
}
int rk(node *o, int val) {
  if (not o) return 0;
  if (o->val < val) {
    return siz(o->le) + 1 + rk(o->ri, val);
  } else {
    return rk(o->le, val);
  }
}
int kth(node *o, int k) {
  if (siz(o->le) < k) return kth(o->ri, k - siz(o->le) - 1);
  if (siz(o->le) == k) return o->val;
  return kth(o->le, k);
}
void insert(node *&o, int val) {
  node *a, *b;
  split(o, a, b, val);
  o = merge(a, merge(new node(val), b));
}
void remove(node *&o, int val) {
  node *a, *b, *c;
  split(o, a, b, val), split2(b, b, c, 1);
  if (b and b->val == val) delete b, b = 0;
  o = merge(a, merge(b, c));
}
node *meld(node *a, node *b) {
  if (not a or not b) return a ? a : b;
  if (rng() % (siz(a) + siz(b)) >= siz(a)) swap(a, b);
  a->push();
  node *c; split(b, b, c, a->val);
  a->le = meld(a->le, b);
  a->ri = meld(a->ri, c);
  a->pull();
  return a;
}
\end{minted}

\newpage

\section{珂朵莉 ODT}

\begin{itemize}[tight_item]
    \item 輸入空區間（$s \ge t$）會 UB，需要先判掉。
\end{itemize}

\begin{minted}{cpp}
struct ODT {
    map<ll, int> mp;
    ODT(ll lim) {  // i in [0, lim)
        mp[0] = -1, mp[lim] = -1;
    }
    auto find(ll i) {
        return prev(mp.upper_bound(i));
    }
    auto &query(ll i) {
        return find(i)->second;
    }
    ll siz(auto it) {
        return next(it)->first - it->first;
    }
    auto split(ll i) {
        return mp.emplace(i, query(i)).first;
    }
    auto make(ll s, ll t) {
        auto tt = split(t);
        auto ss = split(s);
        return make_pair(ss, tt);
    }
    void update(ll s, ll t, auto v) {
        auto [ss, tt] = make(s, t);
        mp.erase(ss, tt);
        mp[s] = v;
    }
    void perform(ll s, ll t) {
        auto [ss, tt] = make(s, t);
        for (auto it = ss; it != tt; ++it) {
            // ...
        }
    }
};
\end{minted}

\section{配對堆}

\begin{minted}{cpp}
struct node {
  ll val; node *ch, *si;
  node(ll v): val(v), ch(0), si(0) {}
};

node *meld(node *a, node *b) {
  if (not a or not b) return a ? a : b;
  if (a == b) return a;
  if (a->val > b->val) swap(a, b);
  b->si = a->ch;
  a->ch = b;
  return a;
}

node *merges(node *a) {
  if (not a or not a->si) return a;
  node *b = a->si, *c = b->si;
  a->si = b->si = 0;
  return meld(merges(c), meld(a, b));
}

ll top(node *o) {
  return o ? o->val : -1;
}

void pop(node *&o) {
  if (not o) return;
  node *t = o;
  o = merges(o->ch);
  delete t;
}
\end{minted}

\section{隨機堆}

\begin{minted}{cpp}
struct node {
  ll val; node *le, *ri;
  node(ll v): val(v), le(0), ri(0) {}
  node *clone() {
    return this;
    // return new node(*this)  // persistent
  }
};

node *merge(node *a, node *b) {
  if (not a or not b) return a ? a : b;
  if (a == b) return a;
  if (a->val > b->val) swap(a, b);
  a = a->clone();
  a->ri = merge(a->ri, b);
  if (rand() & 1) swap(a->le, a->ri);
  return a;
}

ll top(node *o) {
  return o ? o->val : -1;
}

void pop(node *&o) {
  if (not o) return;
  node *t = o;
  o = merge(o->le, o->ri);
  delete t;
}

node *pop(node *o) {  // persistent
  if (not o) return 0;
  node *t = merge(o->le, o->ri);
  delete o;
  return t;
}
\end{minted}

\section{李超樹}

\begin{itemize}[tight_item]
    \item \texttt{tr[N << 2]} 記得開 $4$ 而非 $2$ 倍
    \item \texttt{update} 往下傳 \texttt{v}，而非 \texttt{u}
\end{itemize}

\begin{minted}{cpp}
struct line {
  ll m = 0, k = 0;
  ll y(ll x) {
    return x * m + k;
  }
} tr[N << 2];

void update(int o, int s, int t, line v) {
  auto &u = tr[o];
  int mid = s + t >> 1;
  if (v.y(mid) > u.y(mid)) swap(u, v);
  if (v.y(s+0) > u.y(s+0)) update(o << 1 | 0, s, mid, v);
  if (v.y(t-1) > u.y(t-1)) update(o << 1 | 1, mid, t, v);
}

void modify(int o, int s, int t, int qs, int qt, line v) {
  if (qt <= s or  t <= qs) return;
  if (qs <= s and t <= qt) {
    update(o, s, t, v);
    return;
  }
  int mid = s + t >> 1;
  modify(o << 1 | 0, s, mid, qs, qt, v);
  modify(o << 1 | 1, mid, t, qs, qt, v);
}



ll query(int o, int s, int t, int x) {
  ll res = tr[o].y(x);
  if (s + 1 == t) return res;
  int mid = s + t >> 1;
  if (x < mid) res = max(res, query(o << 1 | 0, s, mid, x));
  else         res = max(res, query(o << 1 | 1, mid, t, x));
  return res;
}
\end{minted}

\section{莫隊}

\begin{itemize}[tight_item]
    \item $B = \frac{N}{\sqrt{Q}} \Rightarrow \mathcal{O}(N \sqrt{Q})$
\end{itemize}

\subsection*{普通莫隊}

\begin{minted}{cpp}
// int B;
struct qq {
    int qs, qt, qi;  // [qs, qt)
    bool operator<(const qq &o) const {
        if (qs / B != o.qs / B) return qs < o.qs;
        if ((qs / B) & 1) return qt < o.qt;
        return qt > o.qt;
    }
} qqs[Q];
ll ans[Q];
// ll cur; void add(int); void sub(int);

sort(qqs, qqs + q);
int s = 0, t = 0;  // [s, t);
for (auto [qs, qt, qi]: qqs) {
    while (s > qs) add(--s);
    while (t < qt) add(t++);
    while (s < qs) sub(s++);
    while (t > qt) sub(--t);
    ans[qi] = cur;
}
\end{minted}

\subsection*{回滾莫隊}

\begin{minted}{cpp}
// int B;
struct qq {
    int qs, qt, qi;  // [qs, qt)
    bool operator<(const qq &o) const {
        if (qs / B != o.qs / B) return qs < o.qs;
        return qt < o.qt;
    }
} qqs[Q];
ll ans[Q];
// ll cur; void add(int); void undo(int);

sort(qqs, qqs + q);
int t = 0;  // [s, t);
int pbi = -1, bt = 0;
for (auto [qs, qt, qi]: qqs) {
    if (int bi = qs / B; bi != pbi) {
        while (t > bt) undo(--t);  // or clear
        pbi = bi, bt = (bi + 1) * B;
        t = bt;
    }
    while (t < qt) add(t++);
    int qtt = min(qt, bt);
    for (int s = qtt; s > qs; ) add(--s);
    ans[qi] = cur;
    for (int s = qs; s < qtt; ) undo(s++);
}
\end{minted}

\newpage

\subsection*{樹上莫隊}

\begin{minted}{cpp}
// int B;
struct qq {
    int qs, qt, x, qi;  // [qs, qt), x = lca
    bool operator<(const qq &o) const {
        if (qs / B != o.qs / B) return qs < o.qs;
        if ((qs / B) & 1) return qt < o.qt;
        return qt > o.qt;
    }
} qqs[Q];
ll ans[Q];
// ll cur; void add(int); void sub(int);

bool vis[N];
// M(vis, 0)
void toggle(int i) {
    if (vis[i] ^= 1) add(i);
    else             sub(i);
}

int tin[N], tout[N], euler[N * 2], dfc;
void dfs(int i, int p) {
    euler[tin[i] = dfc++] = i;
    for (int j: G[i]) if (j != p) dfs(j, i);
    euler[tout[i] = dfc++] = i;
}
// dfc = 0; dfs(0, -1);
REP(qi, 0, q) {  // query a, b
    if (tin[a] > tin[b]) swap(a, b);
    int x = lca(a, b);
    if (x == a) {
        qqs[qi] = {tin[a], tin[b] + 1, -1, qi};
    } else {
        qqs[qi] = {tout[a], tin[b] + 1, x, qi};
    }
}
sort(qqs, qqs + q);
int s = 0, t = 0;  // [s, t);
for (auto [qs, qt, x, qi]: qqs) {
    while (s > qs) toggle(euler[--s]);
    while (t < qt) toggle(euler[t++]);
    while (s < qs) toggle(euler[s++]);
    while (t > qt) toggle(euler[--t]);
    if (~x) toggle(x);
    ans[qi] = cur;
    if (~x) toggle(x);
}
\end{minted}

\section{多重背包 單調隊列}

\begin{minted}{cpp}
const int N = 100, M = 4e4;
int d[M + 1], p[M + 1], q[M];

void solve() {
  int n, m; cin in n in m;
  REP(_, 0, n) {
    int v, w, c; cin in v in w in c;  // 價值 重量 數量
    memcpy(d, p, sizeof(p));
    REP(r, 0, w) {
      for (int i = r, h = 0, t = 0; i <= m; i += w) {
        if (h < t and q[h] < i - w * c) ++h;
        if (h < t) p[i] = max(p[i], d[q[h]] + (i - q[h]) / w * v);
        for (; h < t and d[i] >= d[q[t-1]] + (i - q[t-1]) / w * v; --t);
        q[t++] = i;
      }
    }
  }
  cout ot p[m] nl;  // max sum(v) that sum(w) <= m
}
\end{minted}

\section{Johnson}

\begin{minted}{cpp}
const ll inf = 0x3f3f3f3f'3f3f3f3f;

struct Johnson {
  struct pr {
    int i; ll v;
    bool operator<(const pr &o) const {
      return v > o.v;
    }
  };
  int n;
  vec<vec<pr>> G;
  vec<ll> h;
  vec<vec<ll>> dis;
  Johnson(int nn) : n(nn), G(n), h(n), dis(n) {}
  void add(int i, int j, ll w) {
    G[i].emplace_back(j, w);
  }
  bool run() {
    vec<int> cnt(n);
    deque<int> q;
    vec<bool> inq(n, 1);
    REP(i, 0, n) q.emplace_back(i);
    while (q.size()) {  // SPFA
      int i = q.front(); q.pop_front(); inq[i] = 0;
      ll v = h[i];
      for (auto [j, w]: G[i]) {
        ll nv = v + w;
        if (nv < h[j]) {
          h[j] = nv;
          cnt[j] = cnt[i] + 1;
          if (cnt[j] >= n) return 0;  // neg cyc
          if (not inq[j]) {
            inq[j] = 1;
            q.emplace_back(j);
          }
        }
      }
    }
    priority_queue<pr> pq;
    REP(s, 0, n) {  // Dij * n
      auto &cur = dis[s];
      cur.assign(n, inf);
      cur[s] = 0;
      pq.emplace(s, 0);
      while (pq.size()) {
        auto [i, v] = pq.top(); pq.pop();
        if (v > cur[i]) continue;
        for (auto [j, w]: G[i]) {
          ll nv = v + w + h[i] - h[j];
          if (nv < cur[j]) {
            cur[j] = nv;
            pq.emplace(j, nv);
          }
        }
      }
      REP(i, 0, n) if (cur[i] < inf) cur[i] += h[i] - h[s];
    }
    return 1;  // no neg cycle
  }
};
\end{minted}

\section{MP}

\begin{minted}{cpp}
// s.find(p)
// char p[N], s[N]; int lps[N];
for (int i = 1, j = 0; i < m; ++i) {
  for (; j and p[j] != p[i]; j = lps[j - 1]);
  j += (p[j] == p[i]);
  lps[i] = j;
}
for (int i = 0, j = 0; i < n; ++i) {
  for (; j and p[j] != s[i]; j = lps[j - 1]);
  j += (p[j] == s[i]);
  if (j >= m) {
    ans[a++] = i - (j - 1);
    j = lps[j - 1];
  }
}
\end{minted}

\section{KMP DFA}

\begin{minted}{cpp}
const int A = 26;
array<int, A> nxt[N + 1];

void build(const string &s) {
  int n = s.size();
  nxt[0] = {};
  if (not n) return;
  nxt[0][s[0] - 'a'] = 1;
  int x = 0;
  REP(i, 1, n + 1) {
    nxt[i] = nxt[x];
    if (i < n) {
      char c = s[i] - 'a';
      nxt[i][c] = i + 1;
      x = nxt[x][c];
    }
  }
}
\end{minted}

\section{Z function}

\begin{minted}{cpp}
// char s[N]; int z[N];
fill(z, z + n, 0);
for (int i = 1, l = 0, r = 0; i < n; ++i) {
  if (i < r) z[i] = min(r - i, z[i - l]);
  for (; i + z[i] < n and s[z[i]] == s[i + z[i]]; ++z[i]);
  if (i + z[i] > r) l = i, r = i + z[i];
}
\end{minted}

\section{Minimal Rotation}

\begin{minted}{python}
def min_rotation(s: str) -> str:
  n = len(s)
  s += s
  i, j = 0, 1
  while i < n and j < n:
    k = 0
    while k < n and s[i + k] == s[j + k]: k += 1
    if s[i + k] < s[j + k]: j += k + 1
    else: i += k + 1
    if i == j: j += 1
  i = min(i, j)
  return s[i:i+n]
\end{minted}

\section{馬拉車}

\begin{minted}{cpp}
// char s[N << 1 | 1]; int z[N << 1 | 1];
void pal(char s[], int z[]) {
  int n = strlen(s) << 1 | 1;
  for (int i = n; i--; s[i] = i & 1 ? s[i >> 1] : '#');
  for (int i = 1, l = 0, r = 0; i < n; ++i) {
    if (i < r) z[i] = min(z[2 * l - i], r - i);
    while (i + 1 + z[i] < n and i - 1 - z[i] >= 0 and 
        s[ i + 1 + z[i] ] == s[ i - 1 - z[i]]) ++z[i];
    if (i + z[i] > r) l = i, r = i + z[i];
  }
}
\end{minted}

\section{SCC}

\begin{minted}{cpp}
// int dfn[N], low[N], scc[N], stk[N], dfc, si;
// vec<int> G[N];
// bool instk[N];
// M(dfn, 0);
void tarjan(int i) {
  dfn[i] = low[i] = ++dfc;
  stk[si++] = i, instk[i] = 1;
  for (int j: G[i]) {
    if (not instk[j] and dfn[j]) continue;
    if (not dfn[j]) tarjan(j);
    low[i] = min(low[i], low[j]);
  }
  if (low[i] == dfn[i]) {
    for (int j = -1; j != i; ) {
      j = stk[--si], instk[j] = 0;
      scc[j] = i;
    }
  }
}
\end{minted}

\section{割點}

\begin{minted}{cpp}
// int dfn[N], low[N], dfc;
// vec<int> G[N];
// bool ans[N];
// M(dfn, 0);
void tarjan(int i, int pa) {
  low[i] = dfn[i] = ++dfc;
  int ch = 0;
  for (int j: G[i]) {
    if (not dfn[j]) {
      ++ch;
      tarjan(j, i);
      low[i] = min(low[i], low[j]);
      if (i != pa and low[j] >= dfn[i]) ans[i] = 1;
    } else if (j != pa and dfn[j] < low[i]) {
      low[i] = dfn[j];
    }
  }
  if (i == pa and ch >= 2) ans[i] = 1;
}
\end{minted}

\section{割邊}

\begin{minted}{cpp}
// int dfn[N], low[N], dfc;
// vec<pair<int, int>> G[N];
// bool ans[M];
// M(dfn, 0);
void tarjan(int i, int pa) {
  low[i] = dfn[i] = ++dfc;
  for (auto [j, ei]: G[i]) {
    if (not dfn[j]) {
      tarjan(j, i);
      low[i] = min(low[i], low[j]);
      if (low[j] > dfn[i]) ans[ei] = 1;
    } else if (j != pa and dfn[j] < low[i]) {
      low[i] = dfn[j];
    }
  }
}
\end{minted}

\newpage

\section{點雙}

\begin{minted}{cpp}
// int dfn[N], low[N], stk[N], dfc;
// vec<pair<int, int>> G[N];
// bool instk[N], ans[N];
// M(dfn, 0);
void tarjan(int i, int pa) {
  dfn[i] = low[i] = ++dfc;
  for (auto [j, ei]: G[i]) {
    if (not dfn[j]) {
      stk[si++] = ei;
      tarjan(j, i);
      low[i] = min(low[i], low[j]);
      if (low[j] >= dfn[i]) {
        for (int ej = -1; ej != ei; ++x) {
          ej = stk[--si];
          bcc[ej] = ei;
        }
      }
    } else if (j != pa and dfn[j] < dfn[i]) {
      low[i] = min(low[i], dfn[j]);
      stk[si++] = ei;
    }
  }
}
\end{minted}

\section{邊雙}

\begin{minted}{cpp}
// int dfn[N], low[N], bcc[N], stk[N], dfc, si;
// vec<pair<int, int>> G[N];
// bool instk[N];
// M(dfn, 0);
void tarjan(int i, int inin) {
  dfn[i] = low[i] = ++dfc;
  stk[si++] = i, instk[i] = 1;
  for (auto [j, ei]: G[i]) {
    if (ei == inin) continue;
    if (not instk[j] and dfn[j]) continue;
    if (not dfn[j]) tarjan(j, ei);
    low[i] = min(low[i], low[j]);
  }
  if (low[i] == dfn[i]) {
    for (int j = -1; j != i; ) {
      j = stk[--si], instk[j] = 0;
      bcc[j] = i;
    }
  }
}
\end{minted}

\section{歐拉路徑}

\begin{minted}{python}
for i, it in G:
  G[i] = iter(it)
def rec(i: int) -> None:
  for j in G[i]:
    rec(i)
  path = [i] + path
rec(s)
\end{minted}

\section{tarjan offline DSU}

\begin{minted}{python}
def dfs(i: int, pa: int) -> None:
  top[i] = i
  vis[i] = True
  for qi, j in que[i]:
    if not vis[j]: continue
    ans[qi] = top[find(j)]
  for j in G[i]:
    if j == pa: continue
    dfs(j, i)
  if i is not root:
    ri, rp = find(i), find(pa)
    top[ri] = top[rp]
    merge(ri, rp)
\end{minted}

\section{LCA}

\begin{minted}{cpp}
struct LCA {
  int n;
  vec<int> dep, idx, asc, hd;
  #define lb(x) ((x) & -(x))
  #define hb(x) ((x) ? 1 << __lg(x) : 0)
  LCA(const auto &G, int r = 0): n(G.size()), dep(n), idx(n), asc(n), hd(n + 1) {
    vec<pair<int, int>> qu;
    auto dfs = [&](auto &&self, int i, int p) -> void {
      qu.emplace_back(i, p);
      idx[i] = qu.size();
      for (int j: G[i]) if (j != p) {
        dep[j] = dep[i] + 1;
        self(self, j, i);
        hd[idx[j]] = i;
        if (lb(idx[i]) < lb(idx[j])) idx[i] = idx[j];
      }
    };
    dfs(dfs, r, r);
    for (auto [i, p]: qu) asc[i] = asc[p] | lb(idx[i]);
  }
  int operator()(int a, int b) const {
    if (int x = idx[a] ^ idx[b]) {
      x = asc[a] & asc[b] & -hb(x);
      if (int y = hb(asc[a] ^ x)) y = lb(y), a = hd[(idx[a] & -y) | y];
      if (int y = hb(asc[b] ^ x)) y = lb(y), b = hd[(idx[b] & -y) | y];
    }
    return dep[a] < dep[b] ? a : b;
  }
};
\end{minted}

\section{HLD}

\begin{itemize}[tight_item]
    \item TLE：\texttt{siz[i] += siz[j]}
\end{itemize}

\begin{minted}{cpp}
struct HLD {
  int n, dfc;
  vec<int> pa, ch, top, dep, siz, dfn;
  HLD(const vec<vec<int>> &G) : n(G.size()), dfc(0), pa(n), ch(n, -1), top(n), dep(n), siz(n, 1), dfn(n) {
    auto dfs = [&](auto &&self, int i, int p) -> void {
      pa[i] = p;
      for (int j: G[i]) if (j != p) {
        dep[j] = dep[i] + 1;
        self(self, j, i);
        siz[i] += siz[j];
        if (ch[i] == -1 or siz[j] > siz[ch[i]]) ch[i] = j;
      }
    };
    auto hld = [&](auto &&self, int i, int t) -> void {
      top[i] = t, dfn[i] = dfc++;
      if (~ch[i]) {
        self(self, ch[i], t);
        for (int j: G[i]) if (j != pa[i] and j != ch[i]) {
          self(self, j, j);
        }
      }
    };
    dfs(dfs, 0, -1), hld(hld, 0, 0);
  }


  int lca(int a, int b) {
    while (top[a] != top[b]) {
      if (dep[top[a]] < dep[top[b]]) swap(a, b);
      // [top[a], a]
      a = pa[top[a]];
    }
    if (dep[a] < dep[b]) swap(a, b);
    // [a, b]
    return b;
  }
};
\end{minted}

\section{虛樹}

\begin{minted}{python}
# 關鍵點 -> 虛樹
def build(kp: list[int]) -> defaultdict[list[int]]:
  vt = defaultdict(list)
  kp.sort(key=get_dfn)  # 關鍵點照 dfn 排序
  # 加入 樹根、關鍵點、相鄰關鍵點的 lca
  aux = [root, kp[0]]
  for i in range(1, len(kp)):
    aux.append(kp[i])
    aux.append(lca(kp[i - 1], kp[i]))
  # 虛樹點照 dfn 排序、去重
  aux = sorted(set(aux), key=get_dfn)
  # 對於相鄰 a, b，link(lca(a, b), b)
  for i in range(1, len(aux)):
    x = lca(aux[i - 1], aux[i])
    vt[x].append(aux[i])
  return vt
\end{minted}

\section{圓方樹}

\begin{minted}{cpp}
struct RST {
  int n, dfc, sqr;
  vec<int> dfn, low, stk;
  vec<bool> instk;
  vec<vec<int>> T;
  void link(int a, int b) {
    T[a].emplace_back(b);
    T[b].emplace_back(a);
  }
  RST(const vec<vec<int>> &G) : n(G.size()), dfc(0), sqr(n), dfn(n), low(n), stk(), instk(n), T(n << 1) {
    auto tarjan = [&](auto &&self, int i, int p) -> void {
      dfn[i] = low[i] = ++dfc;
      stk.emplace_back(i);
      for (int j: G[i]) {
        if (not dfn[j]) {
          self(self, j, i);
          low[i] = min(low[i], low[j]);
          if (low[j] >= dfn[i]) {
            for (int x = -1; x != j; ) {
              x = stk.back(); stk.pop_back();
              link(x, sqr);
            }
            link(i, sqr);
            ++sqr;
          }
        } else if (j != p and dfn[j] < dfn[i]) {
          low[i] = min(low[i], dfn[j]);
        }
      }
    };
    REP(i, 0, n) if (not dfn[i]) {
      tarjan(tarjan, i, -1);
      stk.pop_back();
    }
    T.resize(sqr);
  }
};
\end{minted}

\newpage

\section{Boruvka}

\begin{itemize}[tight_item]
    \item 每一輪，每個\textbf{連通塊}（而非單點） 向其他連通塊\textbf{僅}嘗試\textbf{最小的一條邊}，小心不要一次更新兩邊以上。
    \item 連通塊的判斷方式，是在所有 merge 開始之前，所以要另外開一個 array 紀錄 merge 前的所屬 group，而不能使用 dsu find 動態判斷。
    \item 因為是\textbf{相異}連通塊，所以大概率跟維護\textbf{最小、次小}有關，寫法上要注意第一點的「不要一次更新了這兩邊」。
\end{itemize}

\section{Convex Hull}

\begin{minted}{cpp}
struct p {
  int x, y;
  bool operator<(const p &o) const {
    return tie(x, y) < tie(o.x, o.y);
  }
  bool operator==(const p &o) const {
    return tie(x, y) == tie(o.x, o.y);
  }
} l[N];
int h[N << 1];

ll crs(const int o, const int a, const int b) {
  return (ll)(l[a].x - l[o].x) * (l[b].y - l[o].y)
       - (ll)(l[a].y - l[o].y) * (l[b].x - l[o].x);
}

void solve() {
  int n; cin in n;
  REP(i, 0, n) cin in l[i].x in l[i].y;
  sort(l, l + n), n = unique(l, l + n) - l;
  int dn = 0;
  REP(i, 0, n) {
    for (; dn >= 2 and crs(h[dn-2], h[dn-1], i) <= 0; --dn);
    h[dn++] = i;
  }
  int up = dn;
  for (int i = n-1; i--; ) {
    for (; up > dn and crs(h[up-2], h[up-1], i) <= 0; --up);
    h[up++] = i;
  }
  if (up != dn) --up;
}
\end{minted}

\section{Kuhn 二分圖最大匹配}

\begin{minted}{cpp}
const int N = 500;
vec<int> G[N];
int pr[N]; bool vis[N];

bool dfs(int i) {
  for (int j: G[i]) {
    int j = to[ei];
    if (vis[j]) continue;
    vis[j] = 1;
    if (pr[j] == -1 or dfs(pr[j])) {
      pr[j] = i;
      return 1;
    }
  }
  return 0;
}
int kuhn(int n) {
  M(pr, -1);
  int ans = 0;
  REP(i, 0, n) M(vis, 0), ans += dfs(i);
  return ans;
}
\end{minted}

\section{Dinic}

\begin{itemize}[tight_item]
  \item $cap = 1 \ ? \ \mathcal{O}(m \sqrt n) : \mathcal{O}(m n ^ 2)$
\end{itemize}

\begin{minted}{cpp}
const ll inf = 0x3f3f3f3f'3f3f3f3f;

struct Dinic {
  struct ee { int to; ll flow, cap; };
  int n, s, t;
  vec<ee> es;
  vec<vec<int>> G;
  vec<int> dep, it;
  Dinic(int nn) : n(nn), G(n) {}
  bool bfs() {
    dep.assign(n, 0);
    queue<int> q;
    q.emplace(s), dep[s] = 1;
    while (q.size()) {
      int i = q.front(); q.pop();
      for (int ei: G[i]) {
        auto &[j, flow, cap] = es[ei];
        if (dep[j] or flow == cap) continue;
        dep[j] = dep[i] + 1;
        q.emplace(j);
      }
    }
    return dep[t];
  }
  ll dfs(int i, ll lim) {
    if (i == t or not lim) return lim;
    ll res = 0;
    for (int &idx = it[i]; idx < G[i].size(); ++idx) {
      int ei = G[i][idx];
      auto &[j, flow, cap] = es[ei];
      if (dep[j] != dep[i] + 1) continue;
      ll x = dfs(j, min(lim - res, cap - flow));
      res += x;
      es[ei ^ 0].flow += x;
      es[ei ^ 1].flow -= x;
      if (res == lim) break;
    }
    if (res < lim) dep[i] = 0;
    return res;
  }
  void add(int a, int b, ll w) {
    G[a].emplace_back(es.size()), es.emplace_back(b, 0, w);
    G[b].emplace_back(es.size()), es.emplace_back(a, 0, 0);
  }
  ll maxflow(int ss, int tt) {
    s = ss, t = tt;
    ll ans = 0, cur;
    while (bfs()) {
      it.assign(n, 0);
      cur = dfs(s, inf);
      ans += cur;
    }
    return ans;
  }
};
\end{minted}

\section{上下界流}
\begin{minted}{python}
es = [hi - lo for lo, hi in es]  # lo <= flow <= hi
for i, j in es: bal[i] -= lo, bal[j] += lo
if has s-t: aes.add(t, s, inf)  # → .flow = s→t 可行流
for i, b in enum(bal):
  if b > 0: aes.add(ss, i, v), demand += v
  elif v < 0: aes.add(i, tt, -v)
if (可行流 := maxflow(ss, tt, es + aes)) < demand: NO SOL
MAX = 可行流 + maxflow(s, t, es, flow)
MIN = 可行流 - maxflow(t, s, es, flow)
\end{minted}

\section{MCMF}

\begin{minted}{cpp}
const ll inf = 0x3f3f3f3f'3f3f3f3f;

struct MCMF {
  struct ee { int to; ll flow, cap, cost; };
  int n, s, t; ll mc, mf;
  vec<ee> es;
  vec<vec<int>> G;
  vec<ll> h, dis;
  vec<int> it;
  vec<bool> inq;
  MCMF(int nn) : n(nn), G(n), inq(n) {}
  void spfa() {
    queue<int> q;
    q.emplace(s), inq[s] = 1;
    while (q.size()) {
      int i = q.front(); q.pop(), inq[i] = 0;
      for (int ei: G[i]) {
        auto &[j, flow, cap, cost] = es[ei];
        if (flow == cap) continue;
        ll nv = h[i] + cost;
        if (nv >= h[j]) continue;
        h[j] = nv;
        if (not inq[j]) q.emplace(j), inq[j] = 1;
      }
    }
  }
  bool dijkstra() {
    dis.assign(n, inf);
    struct item {
      int i; ll v;
      bool operator<(const item &o) const {
        return v > o.v;
      }
    };
    priority_queue<item> q;
    q.emplace(s, 0), dis[s] = 0;
    while (q.size()) {
      auto [i, v] = q.top(); q.pop();
      if (v > dis[i]) continue;
      for (int ei: G[i]) {
        auto &[j, flow, cap, cost] = es[ei];
        if (flow == cap) continue;
        ll nv = v + cost + h[i] - h[j];
        if (nv >= dis[j]) continue;
        dis[j] = nv;
        q.emplace(j, nv);
      }
    }
    REP(i, 0, n) h[i] += dis[i];
    return dis[t] < inf;
  }
  ll dfs(int i, ll lim) {
    if (i == t or not lim) return lim;
    inq[i] = 1;
    ll res = 0;
    for (int &idx = it[i]; idx < G[i].size(); ++idx) {
      int ei = G[i][idx];
      auto &[j, flow, cap, cost] = es[ei];
      if (inq[j] or h[j] != h[i] + cost) continue;
      ll x = dfs(j, min(lim - res, cap - flow));
      mc += x * cost;
      res += x;
      es[ei ^ 0].flow += x;
      es[ei ^ 1].flow -= x;
      if (res == lim) break;
    }
    inq[i] = 0;
    return res;
  }

  void add(int a, int b, ll w, ll c) {
    G[a].emplace_back(es.size()), es.emplace_back(b, 0, w, +c);
    G[b].emplace_back(es.size()), es.emplace_back(a, 0, 0, -c);
  }
  void maxflow(int ss, int tt) {
    s = ss, t = tt;
    h.assign(n, inf);
    spfa();
    mc = mf = 0;
    while (dijkstra()) {
      it.assign(n, 0);
      ll cur = dfs(s, inf);
      mf += cur;
    }
  }
};
\end{minted}

\section{SPFA NEVER DIES}

\begin{itemize}[tight_item]
    \item \texttt{push} 時與 \texttt{q} 內平均值比較決定頭尾
    \item \texttt{push; dis[q[0]] < dis[q[-1]] (+sqrt sum w) ? swap}
    \item \texttt{2 <= cnt[j] <= sqrt V ? front : back}
\end{itemize}

\section{SA}

\begin{minted}{cpp}
int sa[N], tmp[2][N], c[N], rk[N], lcp[N];
void SA(string &s) {
  int *x = tmp[0], *y = tmp[1], m = 256, n = s.size();
  REP(i, 0, m) c[i] = 0;
  REP(i, 0, n) ++c[x[i] = s[i]];
  REP(i, 1, m) c[i] += c[i - 1];
  for (int i = n; i--; ) sa[--c[x[i]]] = i;
  for (int k = 1; k < n; k <<= 1) {
    REP(i, 0, m) c[i] = 0;
    REP(i, 0, n) ++c[x[i]];
    REP(i, 1, m) c[i] += c[i - 1];
    int p = 0;
    REP(i, n - k, n) y[p++] = i;
    REP(i, 0, n) if (sa[i] >= k) y[p++] = sa[i] - k;
    for (int i = n; i--; ) sa[--c[x[y[i]]]] = y[i];
    y[sa[0]] = p = 0;
    REP(i, 1, n) {
      int a = sa[i], b = sa[i - 1];
      if (not (x[a] == x[b] and a + k < n and b + k < n and x[a + k] == x[b + k])) ++p;
      y[sa[i]] = p;
    }
    if (n == p + 1) break;
    swap(x, y), m = p + 1;
  }
}
void LCP(string &s) {
  // lcp[i] = LCP(sa[i - 1], sa[i])
  // lcp(i, j) = min(lcp[rk[i]+1:rk[j]+1])
  int n = s.size(), val = 0;
  REP(i, 0, n) rk[sa[i]] = i;
  REP(i, 0, n) {
    if (not rk[i]) lca[rk[i]] = 0;
    else {
      if (val) --val;
      int p = sa[rk[i] - 1];
      while (val + i < n and val + p < n and s[val + i] == s[val + p]) ++val;
      lcp[rk[i]] = val;
    }
  }
}
\end{minted}

\newpage

\section{SAM}

\begin{minted}{cpp}
const int A = 26;
struct SAM {
  struct node {
    int len, link, pos; bool vis;
    ll cnt, dp;
    int nxt[A];
    int& operator[](int c) {
      return nxt[c];
    }
  } tr[N << 1];
  int last, siz;

  SAM() {
    M(tr, 0);
    tr[0] = {0, -1, -1, 0, 0};
    last = 0, siz = 1;
  }

  void add(int c) {
    int o = siz++;
    tr[o].len = tr[last].len + 1;
    tr[o].pos = tr[o].len;
    tr[o].cnt = 1;
    int p = last; last = o;
    while (~p and not tr[p][c]) {
      tr[p][c] = o;
      p = tr[p].link;
    }
    if (p == -1) {
      tr[o].link = 0;
    } else {
      int q = tr[p][c];
      if (tr[p].len + 1 == tr[q].len) {
        tr[o].link = q;
      } else {
        int clone = siz++;
        tr[clone] = tr[q];
        tr[clone].len = tr[p].len + 1;
        tr[clone].cnt = 0;
        tr[o].link = tr[q].link = clone;
        while (~p and tr[p][c] == q) {
          tr[p][c] = clone;
          p = tr[p].link;
        }
      }
    }
  }

  void dfs(int o) {
    if (tr[o].vis) return;
    tr[o].vis = 1;
    int pa = tr[o].link, ch;
    tr[o].dp = 1;
    REP(c, 0, A) if (ch = tr[o][c]) {
      dfs(ch);
      tr[o].dp += tr[ch].dp;
    }
    if (~pa) tr[pa].cnt += tr[o].cnt;
  }
} sam;
\end{minted}

\newpage

\section{ACAM}

\begin{itemize}[tight_item]
    \item \texttt{while True}：記得 \texttt{if not o: continue}
\end{itemize}

\begin{minted}{cpp}
// Build Links
// const int FAIL = A, END = A + 1;
// int trie[N+1][A+2], q[N+1];
for (int h = 0, t = 1; h < t; ++h) {
  int o = q[h];
  REP(c, 0, A) {
    if (not trie[o][c]) continue;
    q[t++] = trie[o][c];
    if (not o) continue;  // vital
    int fail = trie[o][FAIL];
    while (not trie[fail][c]) {
      if (not fail) break;
      fail = trie[fail][FAIL];
    }
    if (trie[fail][c]) {
      trie[trie[o][c]][FAIL] = trie[fail][c];
      trie[trie[o][c]][END] |= tr[trie[fail][c]][END];
    }
  }
}
\end{minted}

\section{Digital DP}

\begin{itemize}[tight_item]
    \item \texttt{is\_lim = 1} 時狀態唯一，不用 \texttt{memo}。
\end{itemize}

\begin{minted}{python}
n = ...
num = tuple(map(int, str(n)))

@cache
def dfs(i: int = 0, data: Any = None,
    is_lim: bool = True, start: int = -1,
    debug: tuple[int] = ()) -> int:
  if i == len(num): return 1
  res = 0
  lim = num[i] if is_lim else 9
  for d in range(lim + 1):
    new_data = ...(data, d)
    new_is_lim = is_lim and d == lim
    new_start = start + (start != -1 or d != 0)
    new_debug = debug
    # new_debug += (d,)  # for debugging
    res += dfs(i + 1, new_data, new_is_lim, new_start, new_debug)
  return res
\end{minted}

\section{擴展歐幾里得}

\begin{minted}{cpp}
// ax + by = gcd(a, b)
ll extgcd(ll a, ll b, ll &x, ll &y) {
  x = 1, y = 0;
  ll xx = 0, yy = 1;
  while (b != 0) {
    ll q = a / b;
    a = exchange(b, a % b);
    x = exchange(xx, x - q * xx);
    y = exchange(yy, y - q * yy);
  }
  return a;
}
\end{minted}

\newpage

\section{CRT}

\begin{itemize}[tight_item]
  \item $x = x_i \pmod {m_i}$
\end{itemize}

\begin{minted}{cpp}
pair<ll, ll> CRT(const vec<pair<ll, ll>>& xm) {
  int n = xm.size();
  if (not n) return {0, 1};
  auto mod = [](i128 a, ll b) -> ll {
    return (a % b + b) % b;
  };
  ll x, m;
  REP(i, 0, n) {
    auto [xx, mm] = xm[i];
    xx = mod(xx, mm);
    if (i == 0) { x = xx, m = mm; continue; }
    ll p, q, g = extgcd(m, mm, p, q);
    if ((xx - x) % g) return {-1, -1};  // no sol
    ll s = (xx - x) / g, mg = mm / g;
    ll k = mod((i128)p * s, mg);
    ll lcm = m * mg;
    x = mod(x + (i128)m * k, lcm), m = lcm;
  }
  return {x, m};  // x, lcm(m)
}
\end{minted}

\section{數論分塊}

\begin{minted}{cpp}
// S(n) = sum(f(i) * g(n/i) for i in 1~n)
// let s(n) = sum(f(i) for i in 1~n)
ll s(ll x); ll g(ll x);
ll S(ll n) {
  ll res = 0;
  for (ll le = 1, ri; le <= n; le = ri + 1) {
    ri = n / (n / le);
    res += (s(ri) - s(le - 1)) * g(n / le);
  }
  return res;
}
\end{minted}

\section{乘法反元素}

\begin{minted}{cpp}
inv[1] = 1;
REP(i, 2, lim) inv[i] = (ll)(p - p / i) * inv[p % i] % p;
\end{minted}

\section{Lucas}

\begin{minted}{cpp}
ll lucas(ll n, ll k, ll p) {
  return (k == 0) ? 1 : 
    C(n % p, k % p, p) * lucas(n / p, k / p, p) % p;
}
\end{minted}

\section{Floor Sum}

\begin{itemize}[tight_item]
  \item $\sum^{n-1}_{i=0} \lfloor \frac{a * i + b}{m} \rfloor$ in $\mathcal{O}(n + m + a + b)$
\end{itemize}

\begin{minted}{cpp}
ll floor_sum(ll n, ll m, ll a, ll b) {
  ll ans = 0;
  if (a >= m) ans += (n - 1) * n * (a / m) / 2, a %= m;
  if (b >= m) ans += n * (b / m), b %= m;
  ll ymx = (a * n + b) / m, xmx = (ymx * m - b);
  if (ymx == 0) return ans;
  ans += (n - (xmx + a - 1) / a) * ymx;
  ans += floor_sum(ymx, a, m, (a - xmx % a) % a);
  return ans;
}
\end{minted}

\newpage

\section{質因數分解}

\begin{minted}{cpp}
ll mul(ll x, ll y, ll p) { return (__int128_t)x * y % p; }
vec<ll> chk = {2,325,9375,28178,450775,9780504,1795265022};
ll pw(ll a, ll b, ll n) {
  ll res = 1;
  for (; b; b >>= 1) {
    if (b & 1) res = mul(res, a, n);
    a = mul(a, a, n);
  }
  return res;
}
bool check(ll a, ll d, int s, ll n) {
  a = pw(a, d, n);
  if (a <= 1) return 1;
  REP(i, 0, s) {
    if (a == 1) return 0;
    if (a == n - 1) return 1;
    a = mul(a, a, n);
  }
  return 0;
}
bool isprime(ll n) {
  if (n < 2) return 0;
  if (n % 2 == 0) return n == 2;
  ll d = n - 1, s = 0;
  while (d % 2 == 0) d >>= 1, ++s;
  for (ll i: chk) if (not check(i, d, s, n)) return 0;
  return 1;
}
const vec<ll> smlp = {2, 3, 5, 7, 11, 13, 17, 19};
ll findfac(ll n) {
  if (isprime(n)) return 1;
  for (ll p: smlp) if (n % p == 0) return p;
  ll x, y = 2, d, t = 1;
  for (int l = 2; ; l <<= 1) {
    x = y;
    int m = min(l, 32);
    for (int i = 0; i < l; i += m) {
      d = 1;
      REP(j, 0, m) {
        y = (mul(y, y, n) + t) % n;
        d = mul(d, abs(x - y), n);
      }
      ll g = __gcd(d, n);
      if (g == n) {
        l = 1, y = 2, ++t;
        break;
      }
      if (g != 1) return g;
    }
  }
}
map<ll, int> res;
void rec(ll n) {  // PollardRho
  if (n == 1) return;
  if (isprime(n)) { ++res[n]; return; }
  ll d = findfac(n);
  rec(n / d), rec(d);
}
\end{minted}

\newpage

\section{phi}

\begin{minted}{cpp}
vec<int> ps;  // ps.clear();
bool sieve[lim];  // M(sieve, 1), 
int phi[lim];

phi[1] = 1;
REP(i, 2, lim) {
  if (sieve[i]) {
    ps.emplace_back(i);
    phi[i] = i - 1;
  }
  for (int p: ps) {
    if (i * p >= lim) break;
    sieve[i * p] = 0;
    if (i % p == 0) {
      phi[i * p] = phi[i] * p;
      break;
    }
    phi[i * p] = phi[i] * phi[p];
  }
}
\end{minted}

\section{Cipolla}

\begin{itemize}[tight_item]
  \item $x ^ 2 = n \pmod p$
\end{itemize}

\begin{minted}{cpp}
using i128 = __int128_t;
using pr = pair<ll, ll>;
ll pw(ll aa, ll b, ll mod) {
  ll res = 1;
  i128 a = aa;
  for (a %= mod; b; b >>= 1) {
    if (b & 1) res = res * a % mod;
    a = a * a % mod;
  }
  return res;
}
pr cipolla(ll n, ll p) {
  n = (n % p + p) % p;
  if (n == 0) return {0, 0};
  if (p == 2) return {n, n};
  if (pw(n, p-1 >> 1, p) != 1) return {-1, -1};  // no sol
  static mt19937_64 rng(1145141919810);
  ll a = 0, w = 0;
  do {
    a = rng() % p;
    w = ((i128)a * a - n) % p;
    if (w < 0) w += p;
  } while (pw(w, p-1 >> 1, p) != p - 1);
  auto mul = [&](pr x, pr y) {
    auto [a, b] = x; auto [c, d] = y;
    ll r = ((i128)a * c + (i128)b * d % p * w) % p;
    ll i = ((i128)a * d + (i128)b * c) % p;
    return pr{r, i};
  };
  pr res = {1, 0};
  pr b = {a, 1};
  for (ll x = p / 2 + 1; x; x >>= 1) {
    if (x & 1) res = mul(res, b);
    b = mul(b, b);
  }
  ll x0 = res.first, x1 = (p - x0) % p;
  if (x0 > x1) swap(x0, x1);
  return {x0, x1};
}
\end{minted}

\newpage

\section{Discrete Log}

\begin{itemize}[tight_item]
  \item $a ^ x = b \pmod m$
\end{itemize}

\begin{minted}{cpp}
ll DiscreteLog(ll a, ll b, ll m) {
  ll k = 1 % m, ans = 0, g;
  while ((g = gcd(a, m)) > 1) {
    if (b == k) return ans;
    if (b % g) return -1;
    b /= g, m /= g, ans++, k = (k * a / g) % m;
  }
  if (b == k) return ans;
  const int B = max(1, (int)ceil(sqrt(m)));
  unordered_map<ll, int> vis;
  ll t = 1;
  REP(i, 0, B) vis[t * b % m] = i, t = t * a % m;
  ll cur = k * t % m;
  for (int i = 1; i <= B; ++i, cur = cur * t % m)
    if (vis.count(cur)) return i * B - vis[cur] + ans;
  return -1;
}
\end{minted}

\section{FFT}

\begin{minted}{cpp}
#include <cmath>
#include <complex>
using comp = complex<double>;
const int N = 1 << 20;
const double pi = acos(-1.0);
int revbit[N];
comp a[N], b[N];

void fft(comp *l, int n, int inv) {
  REP(i, 0, n) if (i < revbit[i]) swap(l[i], l[revbit[i]]);
  for (int b = 1; b < n; b <<= 1) {
    comp wn = polar(1.0, inv * pi / b);
    for (int i = 0; i < n; i += b * 2) {
      comp w = 1;
      REP(j, i, i + b) {
        comp x = l[j + 0], y = l[j + b] * w;
        l[j + 0] = x + y;
        l[j + b] = x - y;
        w *= wn;
      }
    }
  }
  if (inv == -1) {
    REP(i, 0, n) l[i] /= n;
  }
}

const int lim = 1 << __lg(n * 2) + 1;
REP(i, 1, lim) revbit[i] = revbit[i >> 1] >> 1 | (i & 1) * (lim >> 1);
fft(l, lim, 1), fft(l, lim, -1)
\end{minted}

\newpage

\section{NTT + CRT}

\begin{minted}{cpp}
using ll = long long;

const int N = 1 << 20;
const ll m1 = 469762049, m2 = 998244353, g = 3;
int revbit[N];
ll a1[N], a2[N], ans[N];

ll pw(ll x, ll n, ll mod) {
  ll res = 1;
  for (x %= mod; n; n >>= 1) {
    if (n & 1) (res *= x) %= mod;
    (x *= x) %= mod;
  }
  return res;
}

void ntt(ll *l, int n, int inv, ll mod) {
  REP(i, 0, n) if (i < revbit[i]) swap(l[i], l[revbit[i]]);
  for (int b = 1; b < n; b <<= 1) {
    ll wn = pw(g, (mod - 1) / (b * 2), mod);
    if (inv == -1) wn = pw(wn, mod - 2, mod);
    for (int i = 0; i < n; i += b * 2) {
      ll w = 1;
      REP(j, i, i + b) {
        ll x = l[j + 0], y = l[j + b] * w % mod;
        l[j + 0] = (x + y) % mod;
        l[j + b] = (x - y + mod) % mod;
        w = w * wn % mod;
      }
    }
  }
  if (inv == -1) {
    ll inv_n = pw(n, mod - 2, mod);
    REP(i, 0, n) l[i] = l[i] * inv_n % mod;
  }
}

const int lim = 1 << __lg(n * 2) + 1;
REP(i, 1, lim) revbit[i] = revbit[i >> 1] >> 1 | (i & 1) * (lim >> 1);

ntt(a1, lim, +1, m1), ntt(a2, lim, +1, m2);
ntt(a1, lim, -1, m1), ntt(a2, lim, -1, m2);

// CRT
ll inv_m1 = pw(m1, m2 - 2, m2);
REP(i, 0, lim) {
  ll k = (a2[i] - a1[i] % m2 + m2) % m2 * inv_m1 % m2;
  ans[i] = a1[i] + k * m1;
}
\end{minted}

\newpage

\section{xor basis}

\begin{minted}{cpp}
template <typename T, int bl>
struct xb {
  array<T, bl> x{};
  int siz = 0;
  void add(T v) {
    if (siz == bl) return;
    for (int i = bl; i--; ) if (v >> i & 1) {
      if (not x[i]) {
        x[i] = v;
        ++siz;
        return;
      }
      v ^= x[i];
    }
  }
  bool find(T v) {
    for (int i = bl; i--; ) if (v >> i & 1) {
      if (not x[i]) return 0;
      v ^= x[i];
    }
    return 1;
  }
  T max() {
    T v = 0;
    for (int i = bl; i--; ) if (not (v >> i & 1)) {
      v ^= x[i];
    }
    return v;
  }
  T kth(T k) {
    T v = 0, cnt = (T)1 << siz;
    for (int i = bl; i--; ) if (x[i]) {
      cnt >>= 1;
      if ((cnt <= k) ^ (v >> i & 1)) v ^= x[i];
      if (cnt <= k) k -= cnt;
    }
    return v;
  }
  T rk(T v) {
    T k = 0, cnt = (T)1 << siz;
    for (int i = bl; i--; ) {
      if (x[i]) cnt >>= 1;
      if (v >> i & 1) {
        k += cnt;
        if (not x[i]) break;
        v ^= x[i];
      }
    }
    return k;
  }
};
\end{minted}

\section{Pick \& Euler}

\begin{itemize}[tight_item]
  \item $\text{面積} = \text{內部格子點} + \dfrac{\text{邊界格子點}}{2} - 1$
  \item $V - E + F = 2$
\end{itemize}

\section{Gosper Hack}

\begin{minted}{cpp}
for (int b = (1 << k) - 1, lim = 1 << n, rr; b < lim;
  rr = b+(b&-b), b = ((rr^b) >> ctz(b)+2)|rr);
\end{minted}

\section{SteinerTree}
\begin{minted}{python}
dp[i][s] = min(dp[i][sub] + dp[i][s ^ sub])
dp[i][s] = min(dp[j][s] + dis[j][i])
\end{minted}

\newpage

\section{定長滑窗}

\begin{minted}{python}
def slide(arr, k, func):
  n = len(arr)
  m = n - k + 1
  ans = []
  for s in range(0, m, k):
    for i in range(s + k - 1, s, -1):
      arr[i - 1] = func(arr[i - 1], arr[i])
    cur = None
    for i in range(s, min(m, s + k)):
      if i > s: cur = func(cur, arr[i + k - 1])
      res = func(arr[i], cur)
      ans.append(res)
  return ans
\end{minted}

\section{實作易錯}

\begin{itemize}[itemsep=0.1em, leftmargin=*]
  \item 多筆測資：先將 \texttt{cin in t} 往下移一行
  \item 多筆測資：提早 \texttt{return} 記得把 input 讀完
  \item 範測沒的特判要手驗、特判記得 \texttt{return}
  \item 範測 WA，確認沒有複製錯。
  \item 輸出任意解：輸出前寫個 checker。
  \item YES / NO 輸出注意格式
  \item 純科技題：最好三人都看過一遍題目，確保理解完全正確
  \item 實作易錯題目：對拍 $\le$ 20 min 則要對拍
  \item 多細節題目：枚舉細節，排列組合亂改，說不定就 AC
  \item 空間 $\approx 10 ^ 8$ 要好好算記憶體
  \item \texttt{continue}、\texttt{break} 容易寫反
  \item 掃描線：注意\textbf{同時刻}下，不同種操作的順序
  \item 做 MST 確保邊權 \texttt{sorted}
  \item 遇到 \textbf{rotate、陣列映射}，讀清楚變數定義
  \item Hash WA：相信 Hash 的正確性，不要相信自己開出 C-Style Array 大小的正確性
  \item 對答案二分搜 $[s, t)$ 區間：\texttt{s + bisect(range(s, t), value, key=...)}，\textbf{前面的 \texttt{s +} 不要忘記}
  \item 求 MEX 先把 $\le n$ 都特判掉
  \item 離散化：
  \begin{itemize}[itemsep=0.1em, leftmargin=*]
    \item 去重？中間被刪掉的空格意義？
    \item \textbf{每個點代表的區間意義！}
    \item 離散化後答案不變，先 \texttt{x = mp(x)} 預處理。
    \item 考慮新增 $\infty$ 作為 sentinel，一個不夠就兩個。
    \item \textbf{戳沒被離散化到的點}，需要考慮戳進\textbf{單位格子}的情況，此時就不用再枚舉前一點。
  \end{itemize}
\end{itemize}

\section{GNOME Settings and Keymap}

\begin{minted}{bash}
gsettings set org.gnome.desktop.
  interface clock-show-seconds true
  wm.keybindings switch-to-workspace-{up/down} []
                 move-to-workspace-{up/down} []
\end{minted}

\begin{itemize}[tight_item]
  \item vscode keymap settings - \texttt{Ctrl K  Ctrl S}
  \item add cursor above/below - \texttt{Ctrl Alt Up/Down}
  \item copy line up/down - \texttt{Shift Alt Up/Down}
\end{itemize}

\end{document}
